论文从选题、资料整理到正文撰写，都得到了指导教师的耐心指导。老师在课题方向把握、I2C
协议理解、FPGA
硬件设计方法以及论文结构安排等方面提出了许多具体建议，使本人能够逐步理清数字系统设计与验证的基本流程，也对毕业设计工作的完整性有了更深入的认识。

还要感谢开源社区对 Verilog I2C interface
项目及相关资料的公开分享。项目中的主机、从机、AXI Lite/Wishbone 封装以及
MyHDL
测试平台，为本文开展模块分析和验证方案整理提供了直接参考。与此同时，同学和朋友在资料查找、环境准备和论文修改阶段给予了许多交流和帮助，在此一并表示感谢。

最后，感谢家人在学习和毕业设计期间给予的支持与鼓励。由于本人能力和时间有限，论文中仍存在不足之处，恳请各位老师批评指正。
